% =========================================================
% Whitepaper: Force Creation via Agentic Autonomy (DoD)
% Source: derived from /home/anboas/AI Agentics.md
% =========================================================
\documentclass[11pt]{article}

\usepackage[letterpaper,margin=1in]{geometry}
\usepackage{microtype}
\usepackage{setspace}
\setstretch{1.12}
\usepackage[T1]{fontenc}
\usepackage[utf8]{inputenc}
\usepackage{lmodern}
\usepackage{xcolor}
\definecolor{BrandAccent}{HTML}{204060}
\definecolor{BrandGray}{HTML}{6E6E6E}
\definecolor{RuleGray}{HTML}{D9D9D9}
\usepackage{fancyhdr}
\pagestyle{fancy}
\fancyhf{}
\renewcommand{\headrulewidth}{0pt}
\fancyfoot[C]{\small\textcolor{BrandGray}{\thepage}}
\usepackage[hidelinks]{hyperref}
\hypersetup{colorlinks=true, linkcolor=BrandAccent, urlcolor=BrandAccent}
\usepackage[most]{tcolorbox}
\newtcolorbox{keytakeaway}{colback=white,colframe=BrandAccent,boxrule=0.8pt,arc=2pt,left=10pt,right=10pt,top=8pt,bottom=8pt}
\newenvironment{execsummary}{\section*{Executive Summary}\vspace{-2pt}}{\vspace{6pt}}

\newcommand{\PaperTitle}{The Unrecognized Shift: From AI Force Multiplication to Force Creation in the Department of Defense}
\newcommand{\PaperSubtitle}{A White Paper on Agentic Autonomy, Trust Scopes, and Strategic Imperatives for Defense}
\newcommand{\AuthorName}{Adam Boas}
\newcommand{\AuthorRole}{Cloud Solutions Architect (defense modernization; agentic governance)}
\newcommand{\PaperDate}{January 2026}

\begin{document}

\thispagestyle{empty}
\vspace*{0.6in}
{\Huge\bfseries\color{BrandAccent}\PaperTitle\par}
\vspace{0.12in}
{\large\color{BrandGray}\PaperSubtitle\par}
\vspace{0.25in}
\noindent\textcolor{RuleGray}{\rule{\textwidth}{0.8pt}}
\vspace{0.2in}
{\normalsize\textbf{\AuthorName}\par\textcolor{BrandGray}{\AuthorRole}\par\vspace{6pt}\textcolor{BrandGray}{\PaperDate}\par}
\vfill
\noindent\textcolor{BrandGray}{\small\textbf{Distribution:} Public.\quad\textbf{Disclaimer:} Views are the author's and do not represent official policy.}
\newpage

\begin{execsummary}
This white paper argues that the Department of Defense is framing AI primarily as \emph{force multiplication} (assistants and efficiency tools) while the strategic inflection point is \emph{force creation}: autonomous, intent-driven agents that execute mission threads within bounded authority and scale primarily with compute.

It proposes a practical Trust Scope Framework (Authority, Consequence, Environment, Uncertainty Tolerance) and an Agent Control Plane that makes autonomy governable: identity/ABAC, tool mediation, policy-as-code guardrails, immutable logs, continuous evaluation, and kill switches.

The policy implication is operational: if the DoD cannot field autonomy with verifiable trust scopes at tempo, adversaries will. The paper provides an adversarial threat model, acquisition guidance (spec-as-deliverable + eval-as-deliverable), and a 12--24 month roadmap aligned to DoD strategic projects.
\end{execsummary}

\section{Introduction}
The DoD is adopting AI, but much of that adoption is anchored in a limited frame: AI as a supervised assistant. That frame captures near-term productivity gains, but it under-prepares the institution for the next regime, where autonomy is the only way to operate inside human relevance windows.

Force creation is not ``bigger copilots.'' It is intent-driven systems that can plan, call tools, and execute actions toward goals with bounded authority.

\section{Paradigm Shift: From ``Build vs Buy'' to ``Write Specs''}
As agentic systems mature, the binding constraint in system creation moves from implementation capacity to specification quality. Humans increasingly provide \emph{intents} and \emph{constraints}; agents generate solutions, test them, and iterate.

In defense contexts, this means shifting investment from manual delivery toward context engineering, evaluation harnesses, and governance layers that keep agent execution inside authorized trust scopes.

\section{Trust Scope Framework}
Trust scope decomposes into four parameters. Each parameter should be measurable and testable:
\begin{itemize}
  \item \textbf{Authority:} what the agent is allowed to do.
  \item \textbf{Consequence:} the impact class of mistakes.
  \item \textbf{Environment:} classification/connectivity/contestation conditions.
  \item \textbf{Uncertainty tolerance:} confidence thresholds and escalation triggers.
\end{itemize}

\begin{keytakeaway}
\textbf{Key takeaway:} Trust is not a vibe. It is a scope with measurable gates.
\end{keytakeaway}

\subsection{Trust Scope Manifest (example)}
\begin{verbatim}
trust_scope_manifest:
  agent_name: "enterprise-agents.v1"
  environment:
    classification: "UNCLAS|CUI|SECRET|TS"
    connectivity: "online|intermittent|offline"
  authority:
    allowed_actions: ["draft", "recommend", "execute_within_budget"]
    prohibited_actions: ["modify_firewall_rules", "release_data_cross_domain"]
  consequence_level: "low|moderate|high|life_safety"
  human_authority:
    approvals_required_for:
      - action: "execute_within_budget"
        threshold: "$10K"
    escalation_triggers:
      - "confidence_below:0.85"
      - "prompt_injection_suspected:true"
  verification:
    acceptance_thresholds:
      task_success_rate: 0.97
      critical_error_rate: 0.001
      audit_log_coverage: 1.0
  controls:
    kill_switch: true
    immutable_audit_logging: true
\end{verbatim}

\section{Agent Control Plane}
To scale autonomy safely, an Agent Control Plane provides:
\begin{itemize}
  \item agent identity + ABAC (least privilege)
  \item tool mediation (allowlists, budgets, sandboxing)
  \item policy-as-code guardrails
  \item immutable audit logs + replay
  \item continuous evaluation + regression gates
  \item kill switch + fallback modes
\end{itemize}

\section{Policy Partitioning}
Partition autonomy into categories with tailored trust scopes:
\begin{itemize}
  \item Category A: enterprise/enabling agents
  \item Category B: intelligence agents
  \item Category C: operational/battle management agents
  \item Category D: weapon system autonomy (special constraints)
\end{itemize}

\section{Adversarial Threat Model}
Key failure modes include prompt injection, tool abuse, data poisoning, supply chain integrity, action forgery, and silent drift. The control plane must detect anomalies, enforce provenance, and manage escalation-rate attacks.

\section{Acquisition: Spec as Deliverable + Eval as Deliverable}
Operationalize ``write specs'' via procurement:
\begin{itemize}
  \item Intent/spec package (constraints, success criteria)
  \item Trust Scope Manifest
  \item Evaluation harness (tests, regressions, red-team scripts)
  \item Transparency artifacts (model/system/data cards)
\end{itemize}

\section{Roadmap (12--24 months)}
\begin{itemize}
  \item Months 1--6: deploy low/moderate consequence trust scopes; build control plane; build eval harness.
  \item Months 7--12: graduate to contested workflows; integrate context engineering.
  \item Months 13--18: scale novel capabilities via competitive mechanisms; expand threat modeling.
  \item Months 19--24: scale data provenance and ATO reciprocity pilots.
\end{itemize}

\section{Policy Recommendations}
\begin{enumerate}
  \item Adopt Trust Scope Manifests as a required artifact per agent family.
  \item Build an Agent Control Plane and make it the default integration surface.
  \item Treat evaluations as a first-class deliverable; block deployment without regression gates.
  \item Run force creation pilots that are explicitly intent-driven and compute-scaled.
\end{enumerate}

\section{Glossary}
\begin{itemize}
  \item \textbf{Force multiplication:} AI increases throughput of human workstreams under constant oversight.
  \item \textbf{Force creation:} AI executes mission-relevant work with delegated authorities, scaling with compute.
  \item \textbf{Trust scope:} Authority \times Consequence \times Environment \times Uncertainty tolerance.
  \item \textbf{Agent control plane:} governance layer for safe scaling.
\end{itemize}

\section{Appendix: Workforce Roles}
\begin{itemize}
  \item Intent/spec engineer
  \item Context engineer
  \item Agent SRE / autonomy operator
  \item AI red team
  \item Assurance lead
\end{itemize}

\section{References}
(Imported from the markdown draft; keep in sync with source.)
\begin{itemize}
  \item DoD AI Strategy (2026). Department of Defense Artificial Intelligence Strategy.
  \item OMB M-26-04 (2025). Advancing governance and risk management for federal AI.
  \item DoD Directive 3000.09 (2023). Autonomy in Weapon Systems.
  \item Newport (2026). Why Didn't AI 'Join the Workforce' in 2025?
  \item VentureBeat (2025). Build vs Buy is dead---AI just killed it.
\end{itemize}

\end{document}
